\hypertarget{appendix-1759f03e8ce7}{}\label{referencestart:appendix-1759f03e8ce7}
\section{Psychophysics and threshold inference}\label{appendix-1759f03e8ce7-psychophysics-and-threshold-inference}

\textbf{Curriculum depth:} Introduced · \textbf{First seen:} XR01 · \textbf{Latest development:} XR01

\subsection{Definition and mechanism}\label{appendix-1759f03e8ce7-definition-and-mechanism}

Psychophysics estimates responses under defined stimuli/tasks. 2AFC chance, lapse, curve and criterion determine a threshold; sensitivity and response bias differ. Nonsignificance is not equivalence, and simulated images cannot independently establish human thresholds.

\subsection{Prerequisites}\label{appendix-1759f03e8ce7-prerequisites}

\begin{itemize}
\tightlist
\item
  \hyperref[appendix-b3aa9e8e01ad]{Contrast sensitivity adaptation and masking} --- Visibility judgments require sensitivity and stimulus-context definitions.
\end{itemize}

\subsection{Used by}\label{appendix-1759f03e8ce7-used-by}

\begin{itemize}
\tightlist
\item
  \hyperref[appendix-8dc848f14052]{Stereoacuity and binocular foveation}
\item
  \hyperref[appendix-65788587ef78]{Research validity and ablation}
\end{itemize}

\subsection{Recall and next depth}\label{appendix-1759f03e8ce7-recall-and-next-depth}

Derive the Weibull-like 75\% threshold and name its assumptions. Pending: approved participant methodology and uncertainty/equivalence analysis. Written coverage is not a claim of learner mastery or accepted numbered-curriculum progress.

\subsection{Source and navigation}\label{appendix-1759f03e8ce7-source-and-navigation}

\hyperref[04-visual-computing-ii-ch042-a-threshold-is-evidence-from-a-procedure-not-a-universal-border]{Expanded teaching explanation} · \hyperref[appendix-fd090d982a15-xr01---independent-seeing-and-representation-source]{Independent coverage} · \href{../01-Bits-Programs-and-Reliable-Systems/main.pdf\#nameddest=appendix-4ad47bdd269a}{Dependency Map (Bits, Programs, and Reliable Systems, Reference notes: Dependency Map)}
